\section{Separation, unfolding, and the center-averaged law}\label{sec:centers}

The pair function in \eqref{eq:full-correlation} has four real arguments.
Each reduction of that information should be stated as an explicit
coordinate change or marginalization. This is especially relevant when
comparing a complex-zero statistic with a real-line spacing law.

\subsection{Center and separation coordinates}

For $u,v\in\C$ set
\begin{equation}\label{eq:center-separation}
 c=\frac{\Re u+\Re v}{2},\quad s=\Re u-\Re v,
 \quad a=\Im u,\quad b=\Im v,
 \qquad
 \mathcal C_{2,\epsilon}^W(c,s,a,b)
    =R_{2,\epsilon}^W(c+s/2+\ii a,c-s/2+\ii b).
\end{equation}
The absolute Jacobian is one. Thus integrating $c$ gives a separation
density with two transverse marks; subsequently integrating $a,b$
gives a horizontal separation density. Neither is the full function.
Dividing any of these by $Z_2>0$ gives its ordered-pair probability
density. The ratio $R_2(u,v)/(r(u)r(v))$, where the denominator is
positive, is a different normalization and does not generally have
total mass one.

A second useful projection integrates the full complex center:
\begin{proposition}[Gaussian separation marginal]\label{prop:separation}
For $s\in\C$,
\begin{equation}\label{eq:separation-gaussian}
 \mathcal D_{2,\epsilon}^W(s)
 :=\int_\C R_{2,\epsilon}^W(c+s/2,c-s/2)\dd A(c)
 =\sum_{j\ne k}w_jw_k
                  G_{\sqrt2\epsilon}(s-(z_j-z_k)).
\end{equation}
Its total mass is $Z_2$.
\end{proposition}
\begin{proof}
For a fixed pair, substitute the two Gaussian factors in the integral.
Completing the square in $c$ gives the density of the difference of
two independent Gaussians, whose covariance is $2\epsilon^2 I$.
Its normalization is $(4\pi\epsilon^2)^{-1}$. Summing and then
integrating in $s$ proves both claims. The map
$(c,s)\mapsto(c+s/2,c-s/2)$ has absolute Jacobian one in $\R^4$.
\end{proof}
Unlike \eqref{eq:center-separation}, this projection integrates out the
common transverse coordinate. It cannot recover two independent marks.

If all zeros under observation are real, write
$g_\epsilon(t)=(2\pi\epsilon^2)^{-1/2}\ee^{-t^2/(2\epsilon^2)}$.
The exact factorization is
\begin{equation}\label{eq:real-zero-correlation}
 R_{2,\epsilon}^W(u,v)
 =g_\epsilon(\Im u)g_\epsilon(\Im v)
   \sum_{j\ne k}w_jw_k
        g_\epsilon(\Re u-x_j)g_\epsilon(\Re v-x_k).
\end{equation}
The transverse spread here is known observation noise; the horizontal
factor retains the full real-zero spacing information. A transversely
vanishing scalar detector would discard this nonzero pair function.

\subsection{A declared center law and an unfolding map}

Let $X$ range over a compact interval $I\subset[10,\infty)$ with a
specified probability measure $\pi$. Examples are normalized Lebesgue
measure on $I$ and a Gaussian density truncated to $I$ and renormalized.
For each $X$, declare the window $W_X$, smoothing width $\epsilon_X>0$,
and positive horizontal and transverse scales $\ell_X,b_X$. Define
\begin{equation}\label{eq:unfolding}
 T_X(\zeta)=X+\ell_X\Re\zeta+\ii b_X\Im\zeta,
 \quad J_X=\ell_X b_X,\qquad
 \widetilde R_{k,X}(\boldsymbol\zeta)
   =J_X^kR_{k,\epsilon_X}^{W_X}(T_X\zeta_1,\ldots,T_X\zeta_k).
\end{equation}
The Jacobian preserves total factorial mass. It is legitimate to choose
$\ell_X$ as a nominal mean-spacing scale, but that choice does not
prove that the actual local mean density is constant. We smooth first
in physical coordinates, then unfold; the two resulting Gaussian widths
are $\epsilon_X/\ell_X$ and $\epsilon_X/b_X$. Fixed isotropic smoothing
after unfolding is a different convention, implemented by an anisotropic
physical Gaussian.

Average the conditional observation laws over $X$. Their correlation
functions, including all orders, are
\begin{equation}\label{eq:center-average}
 \overline R_k(\boldsymbol\zeta)
       =\int_I\widetilde R_{k,X}(\boldsymbol\zeta)\dd\pi(X).
\end{equation}
This is a fully specified statistical ensemble derived from the
canonical configuration and an external center law. For simplicity,
all parameter choices below are measurable and bounded on $I$, with
positive lower bounds on the scales and smoothing width. A Gaussian
law on the whole real line would require an additional center-tail
bound; truncation is part of the definition here.

\begin{proposition}[The connected center-averaged pair]\label{prop:covariance}
Let $\widetilde r_X(\zeta)=J_Xr_{\epsilon_X}^{W_X}(T_X\zeta)$ and
$\widetilde d_X(\zeta,\eta)=J_X^2d_{\epsilon_X}^{W_X}(T_X\zeta,T_X\eta)$.
Then, with $\overline r=\int_I\widetilde r_X\dd\pi$,
\begin{equation}\label{eq:center-covariance}
 \overline R_2(\zeta,\eta)-\overline r(\zeta)\overline r(\eta)
 =\operatorname{Cov}_{\pi}\bigl(\widetilde r_X(\zeta),
                              \widetilde r_X(\eta)\bigr)
       -\int_I\widetilde d_X(\zeta,\eta)\dd\pi(X).
\end{equation}
Here covariance is the bilinear real covariance of the displayed
real-valued densities.
\end{proposition}
\begin{proof}
Multiply \eqref{eq:pair-rd} by $J_X^2$, substitute $T_X\zeta,T_X\eta$,
and average. Subtract the product of the two averaged first intensities.
The remaining product term is exactly the covariance in
\eqref{eq:center-covariance}.
\end{proof}

At a fixed center the connected function for the auxiliary independent
observations is just $-d$. The full pair function still contains the
configuration's spatial geometry. Averaging over centers adds the
covariance of that geometry. Thus a canonical estimate of the
center covariance is a concrete further analytical target; it is not
provided by independence of the observation noise. A source expression
for this covariance is obtained in Section~\ref{sec:full-pair}.

The CUE calculation in Appendix~\ref{app:ensemble} provides an explicit
comparison law derived from a supplied ensemble. No such ensemble law
is assumed for the canonical divisor. Its center covariance requires
the additional information isolated in Section~\ref{sec:joint-source}.
